\documentclass[11pt]{amsart}
\usepackage{amssymb}
\usepackage[margin=1.2in]{geometry}

\newtheorem{theorem}{Theorem}
\newtheorem{lemma}[theorem]{Lemma}
\newtheorem{corollary}[theorem]{Corollary}
\theoremstyle{remark}
\newtheorem{remark}[theorem]{Remark}

\title{Rank-one Horrocks over any commutative ring,\\ by a local-global principle for comaximal monoids}
\author{Thierry Coquand}
\date{}

% stamp on every page: file name and date/time of compilation
\newcount\stamphour \newcount\stampmin
\stamphour=\time \divide\stamphour by 60
\stampmin=\time \multiply\stamphour by 60 \advance\stampmin by -\stamphour \divide\stamphour by 60
\edef\stamptime{\the\year-\ifnum\month<10 0\fi\the\month-\ifnum\day<10 0\fi\the\day\space
 \ifnum\stamphour<10 0\fi\the\stamphour:\ifnum\stampmin<10 0\fi\the\stampmin}
\AddToHook{shipout/foreground}{\put(0.5\paperwidth,-\dimexpr\paperheight-0.45in\relax)%
 {\makebox[0pt]{\scriptsize\ttfamily\jobname.pdf\quad\stamptime}}}
\begin{document}

\begin{abstract}
We give a constructive proof that an ideal of $A[X]$ dividing a principal ideal
generated by a monic polynomial is principal, for an arbitrary commutative
ring $A$. The local part follows Nashier--Nichols. The global part uses the
comaximal monoids produced by a formal gcd computation directly, through two
general lemmas: membership in a finitely generated ideal is local, and
compatible local elements glue. The gluing lemma has an explicit proof without
principal localisations and without exponents. This organisation of the proof
was suggested by writing a complete machine-checked version.
\end{abstract}

\maketitle

\section{Statement}

Let $A$ be a commutative ring. All rings are commutative.

\begin{theorem}\label{thm:main}
Let $F\in A[X]$ be monic, and $y_1,\dots,y_n$, $x_1,\dots,x_n\in A[X]$ such that
\[
\textstyle\sum_i x_iy_i = F,\qquad x_iy_j = F e_{ij}\quad(1\le i,j\le n)
\]
for some $e_{ij}\in A[X]$. Then the ideal $L=(y_1,\dots,y_n)$ is principal.
Moreover there are $a_i,b_i\in A[X]$ with $b_ia_j=e_{ij}$ and $\sum_i a_ib_i=1$.
\end{theorem}

The hypothesis says $L\cdot M=(F)$ for $M=(x_1,\dots,x_n)$. If $e$ is an
idempotent matrix of rank one over $A[X]$ which becomes free over $A[X]_f$,
with $f$ monic, then there are $x,y$ satisfying the hypothesis with $F=f^k$,
and the conclusion says that $e$ is free. This is the rank-one case of
Horrocks' theorem used in \cite{Proj}; the projector equations themselves are
not needed. The generator of $L$ need not be monic: if $\varepsilon^2=\varepsilon$
then $(\varepsilon X+1-\varepsilon)$ divides $(X)$.

We write $S^{-1}A$ or $A_S$ for the localisation at a multiplicative monoid $S$.
Recall that $a/1=0$ in $A_S$ means $sa=0$ for some $s\in S$.

\section{Comaximal monoids}

For finite sets $I,U\subseteq A$, let $M(U)$ be the monoid generated by $U$ and
\[
S(I;U) = M(U) + (I).
\]
This is a monoid, and $S(I;U)+(I)=S(I;U)$, so the ideal generated by $I$ in
$A_{S(I;U)}$ is contained in its Jacobson radical: for $c\in(I)$ and $s\in S$,
$1+c/s=(s+c)/s$ is invertible.

A \emph{tree} is a finite binary tree whose inner nodes are labelled by elements
of $A$. Its nodes carry pairs $(I;U)$: the root carries $(\emptyset;\emptyset)$
and a node $(I;U)$ labelled $a$ has the children $(I,a;U)$ and $(I;U,a)$.

\begin{lemma}[{\cite[Chapter XV]{LQ15}}]\label{lem:comax}
Choose, for each leaf $(I;U)$ of a tree, an element $s$ of $S(I;U)$. Then the
chosen elements generate the unit ideal.
\end{lemma}

\begin{proof}
Let $\mathfrak a$ be the ideal generated by the chosen elements. We show by
induction on the tree that $\mathfrak a$ meets $S(I;U)$ for every node; at the
root $S(\emptyset;\emptyset)=\{1\}$. It suffices to show: if $s\in S(I,a;U)$ and
$t\in S(I;U,a)$, then $(s,t)$ meets $S(I;U)$. Write $s=u+i+ab$ and
$t=u'a^k+i'$ with $u,u'\in M(U)$ and $i,i'\in(I)$. Expanding
$u^k=(s-i-ab)^k$ gives $u^k=cs+j+(-b)^ka^k$ with $j\in(I)$, hence
\[
cu's+(-b)^kt = u'u^k + (-u'j + (-b)^ki') \in S(I;U).\qedhere
\]
\end{proof}

The proof never decides anything about elements of $A$. A finite family of
monoids with this property is called \emph{comaximal}.

\section{Local generators}

In this section $R$ is a ring, $J$ an ideal contained in its Jacobson radical,
and we write $\bar{\ }$ for reduction modulo $J$. A monic polynomial is regular,
also modulo $J$.

\begin{lemma}[Monic lifting {\cite[IV.5.1]{Lam06}, \cite[Lemma 7.5]{Proj}}]\label{lem:lift}
Let $L\subseteq R[X]$ be an ideal containing a monic polynomial. If
$\rho\in L$ has all coefficients of degree $>r$ in $J$ and its coefficient of
degree $r$ invertible modulo $J$, then $L$ contains a monic polynomial $h$ of
degree $r$ with $\bar h=\bar\rho\,\bar c$ for some unit $c$.
\end{lemma}

\begin{lemma}[{\cite[IV.5.2]{Lam06}, \cite[Lemma 7.6]{Proj}}]\label{lem:local}
Let $y,x,e,F$ over $R[X]$ satisfy the hypothesis of Theorem~\ref{thm:main}, and
let $h=\sum_jc_jy_j$ be monic with $y_i-ha_i\in J[X]$ for all $i$. Then
$y_i\in(h)$ for all $i$, so $L=(h)$.
\end{lemma}

\begin{proof}
First, if $p=\sum_j c'_jy_j$ lies in $J[X]$ then $p\in J[X]\,L$. Indeed put
$q_i=\sum_je_{ij}c'_j$, so that $x_ip=Fq_i$. Then $\bar F\bar q_i=0$, hence
$q_i\in J[X]$, and $Fp=\sum_iy_ix_ip=F\sum_iq_iy_i$ gives $p=\sum_iq_iy_i$.
Applied to $p=y_i-ha_i$ this gives
\[
y_i = ha_i + \textstyle\sum_jq_{ij}y_j,\qquad q_{ij}\in J[X].
\]
Let $d=\deg h$ and let $v_i$ be the image of $y_i$ in $R[X]/(h)$. The
$R$-module $N$ generated by the $X^kv_i$, $k<d$, is closed under
multiplication by $X$, since $h$ is monic. From $v_i=\sum_jq_{ij}v_j$ we get
$N=JN$, and $N$ is finitely generated, so $N=0$ by Nakayama's lemma. Hence
$y_i\in(h)$.
\end{proof}

\begin{lemma}\label{lem:unique}
If $h,h'\in R[X]$ are monic of degrees $d,d'$ and $(h)=(h')$, then $h=h'$ if
$d=d'$, and $R=0$ if $d\ne d'$.
\end{lemma}

\begin{proof}
A nonzero multiple of a monic polynomial of degree $d$ has degree at least $d$:
comparing coefficients from the top shows that a multiple of degree $<d$ is $0$.
Write $h=h'a$. If $d<d'$ (the case $d>d'$ is symmetric), then $h=0$ by this remark, and $1=0$ since $h$ is
monic. If $d=d'$, write also $h'=hb$; then $h-h'=h'(a-1)$ has degree $<d'$,
so $h=h'$.
\end{proof}

\section{Formal gcd}

We keep the hypothesis of Theorem~\ref{thm:main}. For a node $(I;U)$ write
$S=S(I;U)$ and $J=(I)A_S$.

\begin{lemma}\label{lem:tree}
There is a tree such that at every leaf $(I;U)$ there is a monic $h\in L A_S[X]$
with $y_i-ha_i\in J[X]$ for some $a_i\in A_S[X]$.
\end{lemma}

\begin{proof}
At a node we are given a monic $g\in LA_S[X]$; at the root $g=F$. Divide each
$y_i$ by $g$, with remainders of degree $<\deg g$, and inspect their
coefficients from the top. Each coefficient is a fraction $a/s$, and we branch
on $a$: in the child $(I,a;U)$ it lies in $J$, in the child $(I;U,a)$ it is
invertible. If all coefficients lie in $J$ we have a leaf with $h=g$. Otherwise
some remainder has its highest coefficient not in $J$ invertible, at degree
$r<\deg g$, and Lemma~\ref{lem:lift} gives a monic member of degree $r$, with
which we start again. The degree decreases, so this terminates.
\end{proof}

By Lemma~\ref{lem:local}, applied over $R=A_S$, we get at each leaf a monic
$h$ with $(h)=LA_S[X]$.

\section{A local-global principle}

In this section $S_1,\dots,S_m$ are comaximal monoids of $A$: any choice of
$s_i\in S_i$ generates the unit ideal. Let $B$ be an $A$-algebra and write
$B_i=S_i^{-1}B$ and $B_{ij}=(S_iS_j)^{-1}B$, where $S_iS_j$ is the monoid of
products. We use two facts about $B_i$: finitely many elements of $B_i$ can be
written $p_1/d,\dots,p_r/d$ with one $d\in S_i$, and $p/1=0$ in $B_i$ means $sp=0$
for some $s\in S_i$.

\begin{lemma}\label{lem:member}
Let $z,w_1,\dots,w_r\in B$. If $z\in(w_1,\dots,w_r)B_i$ for every $i$, then
$z\in(w_1,\dots,w_r)$.
\end{lemma}

\begin{proof}
Fix $i$. Clearing denominators in $z=\sum_ja_jw_j$ over $B_i$ gives $d\in S_i$
and $p_j\in B$ with $dz-\sum_jp_jw_j=0$ in $B_i$; this difference is killed by
some $t\in S_i$. Hence $s_iz=\sum_ja_{ij}w_j$ in $B$, with $s_i=td\in S_i$ and
$a_{ij}=tp_j$. Write $\sum_i\lambda_is_i=1$; then
$z=\sum_j(\sum_i\lambda_ia_{ij})w_j$.
\end{proof}

\begin{lemma}\label{lem:glue}
Let $P_i\in B_i$ be such that $P_i$ and $P_j$ have the same image in $B_{ij}$
for all $i,j$. Then there is $g\in B$ with $g=P_i$ in $B_i$ for every $i$.
\end{lemma}

\begin{proof}
Write $P_i=p_i/d_i$ with $p_i\in B$ and $d_i\in S_i$. In $B_{ij}$ we have
$d_jp_i-d_ip_j=d_id_j(P_i-P_j)=0$, so there are $u_{ij}\in S_i$ and
$v_{ij}\in S_j$ with
\[
u_{ij}v_{ij}\,(d_jp_i-d_ip_j)=0\quad\text{in }B.
\]
Choose such a pair for every ordered pair $(i,j)$ and put
\[
r_i=\textstyle\prod_k u_{ik}v_{ki}\in S_i,\qquad s_i=d_ir_i\in S_i,\qquad H_i=r_ip_i .
\]
Since $u_{ij}$ divides $r_i$ and $v_{ij}$ divides $r_j$, the product $r_ir_j$ is
a multiple of $u_{ij}v_{ij}$, so $r_ir_j(d_jp_i-d_ip_j)=0$, that is
$s_iH_j=s_jH_i$ in $B$. Write $\sum_i\lambda_is_i=1$ and put
$g=\sum_j\lambda_jH_j$. Then
\[
s_ig=\textstyle\sum_j\lambda_js_jH_i=H_i=r_ip_i,
\]
and since $s_i=d_ir_i$ is invertible in $B_i$, $g=p_i/d_i=P_i$ in $B_i$.
\end{proof}

Both lemmas say that $B\to\prod_iB_i$ behaves like the restriction to a cover.
The point of Lemma~\ref{lem:glue} is that it uses one element of each monoid
for each pair, and no common exponent.

\section{Proof of Theorem~\ref{thm:main}}

Take the tree of Lemma~\ref{lem:tree}, number its leaves $1,\dots,m$ and let
$S_1,\dots,S_m$ be their monoids; they are comaximal by Lemma~\ref{lem:comax}.
Let $B=A[X]$, so $B_i=A_{S_i}[X]$. At leaf $i$ we have a monic generator $h_i$
of $LB_i$.

\emph{The local generators agree.} Over $B_{ij}=A_{S_iS_j}[X]$, both $h_i$ and
$h_j$ are monic generators of $LB_{ij}$, so they are equal by
Lemma~\ref{lem:unique}; if their degrees differ, then $B_{ij}=0$ and they are
equal anyway.

\emph{Gluing.} By Lemma~\ref{lem:glue} there is $g\in A[X]$ with $g=h_i$ in
every $B_i$. Then $g\in LB_i$ and $y_j\in(g)B_i$ for all $i,j$, so
Lemma~\ref{lem:member}, applied twice, gives $g\in L$ and $y_j\in(g)$. Hence
$L=(g)$.

\emph{Descent.} Write $g=\sum_jc_jy_j$ and $y_j=ga_j$, and put
$b_i=\sum_je_{ij}c_j$, so that $x_ig=Fb_i$. Then $Fb_ia_j=x_iy_j=Fe_{ij}$ and
$F\sum_ia_ib_i=\sum_ix_iy_i=F$; since $F$ is regular, $b_ia_j=e_{ij}$ and
$\sum_ia_ib_i=1$. This proves Theorem~\ref{thm:main}.

\section{Remarks}

The usual way to finish (see \cite[Corollary 7.7]{Proj}) is to choose one element
$s_i$ in each monoid such that the generators descend to $A_{s_i}[X]$, and then
to patch over the principal localisations, raising all elements to a common
power. The organisation above was suggested by the formalisation of the proof in
a small dependent type theory. There $A_S$ is a quotient of pairs, and an
equation in $A_S$ comes with one element of $S$ as witness. So the data at hand
are finitely many elements of each monoid, and it is simplest to multiply them
inside the monoid. Stated for any $A$-algebra, Lemmas~\ref{lem:member}
and~\ref{lem:glue} reduce the proof of the theorem to a short sequence of
steps, and are reusable. The formal statements only use the two facts about $B_i$ listed above, so they
apply to any family of maps $B\to D_i$ with these properties.

The formalisation also showed that only the equations of
Theorem~\ref{thm:main} are used, not that $e$ is a projector, and that
Lemma~\ref{lem:unique} must allow different degrees, since equality of the
degrees of the local generators cannot be decided. The complete proof is
checked, from the definition of polynomial rings and localisations, in about
$5\,000$ lines, in under $3$ seconds.

\begin{thebibliography}{9}
\bibitem{Proj} F.~Cherubini, T.~Coquand, M.~Ritter, D.~W\"arn.
\emph{Projective space in synthetic algebraic geometry}. 2026.
\bibitem{Lam06} T.~Y.~Lam. \emph{Serre's Problem on Projective Modules}.
Springer, 2006.
\bibitem{LQ15} H.~Lombardi, C.~Quitt\'e. \emph{Commutative Algebra:
Constructive Methods}. Springer, 2015.
\end{thebibliography}

\end{document}
